\chapter{Competencias desarrolladas}

Durante la estadía como residente profesional en la empresa TenarisTamsa se adquirieron experiencias en el fortalecimiento de las capacidades profesionales, permitiendo una integración directa con el entorno laboral y proporcionando una visión más amplia sobre los procesos y dinámicas presentes dentro de la industria tecnológica y empresarial. La participación en proyectos reales y en actividades colaborativas dentro del área de Operational Planning IT contribuyó de manera importante al desarrollo tanto técnico como personal, facilitando la adaptación a un entorno profesional enfocado en la resolución de problemas y mejora continua.

Durante el desarrollo del proyecto se aplicaron diversas competencias técnicas adquiridas a lo largo de la formación académica, fortaleciendo conocimientos relacionados con programación, desarrollo de aplicaciones web y estructuración de sistemas de software. Entre las principales habilidades técnicas utilizadas destacan el uso de lenguajes de programación orientados al desarrollo front-end y back-end, la implementación de estructuras de datos, la aplicación de principios de programación orientada a objetos y el manejo de herramientas modernas para la construcción de aplicaciones web.

Asi como el reforzamiento del conocimiento en metodologías ágiles de desarrollo, comprendiendo la importancia de la organización de tareas, seguimiento de actividades y colaboración continua entre los integrantes del equipo para cumplir con los objetivos establecidos dentro de cada fase del proyecto. Esto permitió desarrollar una mayor capacidad de adaptación ante cambios en requerimientos y necesidades operativas del sistema.

De igual manera, el proyecto permitió fortalecer competencias relacionadas con el análisis y resolución de problemas, ya que durante el proceso de desarrollo fue necesario identificar errores, optimizar procesos y proponer soluciones funcionales alineadas con las necesidades del área. Estas actividades favorecieron el pensamiento lógico, la capacidad de investigación y la toma de decisiones técnicas dentro de un entorno de trabajo real.

Por otro lado, también se desarrollaron habilidades blandas fundamentales para el desempeño profesional, tales como la comunicación efectiva, el trabajo en equipo y la colaboración interdisciplinaria. La interacción constante con compañeros del área y con otras áreas involucradas en el desarrollo del sistema permitió mejorar la capacidad para compartir ideas, recibir retroalimentación y coordinar actividades de manera eficiente dentro de un ambiente colaborativo.